% =====================================================================
\section{Discussion}
% =====================================================================

\subsection{Threats to Validity}
The authors are candid about the limits of what the case study can
show, and the study identifies four major threats to validity.

\emph{Lack of a controlled comparison.} The study is not a controlled
experiment against an alternative migration technique, so the
individual contribution of each pipeline component, reference
discovery, categorization, and the LLM itself, cannot be cleanly
isolated from one another.

\emph{Narrow scope.} The study is confined to a single organization's
internal tooling, a single internally fine-tuned model, and a single
kind of migration (widening an identifier's bit width). The authors
caution that results may not transfer to other migration types, such
as framework upgrades or dependency updates, which touch code in
qualitatively different ways.

\emph{Participant and developer bias.} Only three developers took
part, and all three were also involved in building the system itself,
which could make their satisfaction ratings more favorable than an
arms-length evaluation would produce.

\emph{Imprecise effort measurement.} The headline ``50\% time
savings'' figure is based on developer perception rather than measured
wall-clock time. The paper's chosen proxy for effort, Levenshtein edit
distance between file snapshots, is easy to compute automatically but
arguably under- or over-counts the true cognitive load of reviewing,
understanding, and fixing a change.

\subsection{Generalizability and Future Directions}
Despite these caveats, the reported workflow, cheap, imprecise
reference discovery paired with an LLM that fills the resulting
precision gap, wrapped in an inexpensive-to-expensive chain of
automated validators, may provide a reusable framework for other
large-scale migration tasks rather than a solution specific to
32-to-64-bit identifier migrations. The paper's own suggested future
directions reinforce this framing: larger context windows, better
support for lower-resource languages such as Dart, techniques to
detect and suppress hallucinated no-op edits, and, most ambitiously,
using an LLM not just to \emph{make} the edits but also to help
\emph{decide where} edits are needed, which could shrink the
left-over bucket that still requires manual triage. More broadly, the
case study suggests that for large-scale code maintenance, the
highest-leverage design choice may not be maximizing the precision of
the static analysis used to locate candidate edits, but rather
accepting an intentionally imprecise, high-recall search and letting
a general-purpose LLM, backed by cheap automated checks, absorb the
resulting precision gap.